% ===========================================================================
% section_cau1.tex — Câu 1: Multi-Task Vietnamese NLP with PhoBERT
% ===========================================================================
% PACKAGES CẦN THÊM VÀO main.tex (preamble):
% \usepackage{amsmath, amssymb}
% \usepackage{graphicx}
% \usepackage{booktabs}
% \usepackage{multirow}
% \usepackage{hyperref}
% \usepackage{tikz}
% \usetikzlibrary{positioning, arrows.meta, shapes.geometric, fit, calc}
% \usepackage{listings}
% \usepackage{xcolor}
% \usepackage{subcaption}
% \usepackage{float}
% ===========================================================================

\section{Câu 1: Multi-Task Learning cho NLP Tiếng Việt với PhoBERT}
\label{sec:cau1}

\subsection{Giới thiệu bài toán}

Bài toán yêu cầu xây dựng một hệ thống NLP đa nhiệm vụ (multi-task) cho tiếng Việt,
sử dụng chung một backbone PhoBERT và ba đầu ra (head) chuyên biệt cho ba tác vụ:

\begin{enumerate}
    \item \textbf{Named Entity Recognition (NER)}: Nhận diện thực thể có tên (PER, ORG, LOC, MISC) trên tập VLSP 2016 NER theo lược đồ IOB2.
    \item \textbf{Sentiment Analysis}: Phân loại cảm xúc (Negative, Neutral, Positive) trên tập UIT-VSFC.
    \item \textbf{Aspect-Based Sentiment Analysis (ABSA)}: Phân tích cảm xúc theo 10 khía cạnh sản phẩm trên tập UIT-ViSFD.
\end{enumerate}

Ý tưởng chính là \emph{parameter sharing}: backbone PhoBERT được chia sẻ giữa ba tác vụ,
giúp mô hình học được biểu diễn ngôn ngữ chung phong phú hơn so với huấn luyện riêng lẻ,
đồng thời giảm tổng số tham số cần lưu trữ.

% -------------------------------------------------------------------
\subsection{Cơ sở lý thuyết}
% -------------------------------------------------------------------

\subsubsection{PhoBERT — Pre-trained Language Model cho Tiếng Việt}

PhoBERT \cite{phobert2020} là mô hình ngôn ngữ dựa trên kiến trúc RoBERTa,
được huấn luyện trước trên 20\,GB văn bản tiếng Việt. Mô hình sử dụng:
\begin{itemize}
    \item Kiến trúc Transformer Encoder 12 lớp (base) với $H = 768$.
    \item Tokenizer dựa trên BPE (Byte-Pair Encoding) với từ điển 64K subword.
    \item Huấn luyện trước bằng Masked Language Modeling (MLM).
\end{itemize}

Biểu diễn ngữ cảnh cho token thứ $i$ được tính qua cơ chế Multi-Head Self-Attention:
\begin{equation}
    \text{Attention}(Q, K, V) = \text{softmax}\!\left(\frac{QK^\top}{\sqrt{d_k}}\right) V
\end{equation}
trong đó $Q = XW^Q$, $K = XW^K$, $V = XW^V$ với $d_k = H / n_{\text{heads}}$.

\subsubsection{Multi-Task Learning}

Multi-task learning tối ưu hàm mục tiêu kết hợp:
\begin{equation}
    \mathcal{L}_{\text{total}} = \sum_{t \in \{\text{ner}, \text{sent}, \text{absa}\}} w_t \cdot \mathcal{L}_t
\end{equation}
với các trọng số lấy mẫu $w_{\text{ner}} = 0.3$, $w_{\text{sent}} = 0.4$, $w_{\text{absa}} = 0.3$
thực hiện qua \emph{alternating batch sampling}.

\subsubsection{Conditional Random Field (CRF) cho NER}

Đối với NER, thay vì phân loại độc lập từng token, ta sử dụng CRF
để mô hình hoá quan hệ chuyển tiếp giữa nhãn liền kề.
Xác suất của chuỗi nhãn $\mathbf{y}$ cho chuỗi quan sát $\mathbf{x}$:
\begin{equation}
    P(\mathbf{y} \mid \mathbf{x}) = \frac{
        \exp\!\bigl(\sum_{i=1}^{n} \phi(y_i, \mathbf{x}, i) + \sum_{i=1}^{n-1} T_{y_i, y_{i+1}}\bigr)
    }{
        \sum_{\mathbf{y}'} \exp\!\bigl(\sum_{i=1}^{n} \phi(y'_i, \mathbf{x}, i) + \sum_{i=1}^{n-1} T_{y'_i, y'_{i+1}}\bigr)
    }
\end{equation}
trong đó $\phi(y_i, \mathbf{x}, i)$ là emission score từ lớp Linear, còn $T$ là ma trận chuyển tiếp được học.
Giải mã bằng thuật toán Viterbi với độ phức tạp $O(nC^2)$, $C$ là số nhãn.

\subsubsection{Chiến lược Pooling cho Sentiment và ABSA}

\begin{itemize}
    \item \textbf{Sentiment}: Mean-pooling có che (masked) qua tất cả token không phải padding:
    \begin{equation}
        \mathbf{h}_{\text{sent}} = \frac{\sum_{i=1}^{L} m_i \cdot \mathbf{h}_i}{\sum_{i=1}^{L} m_i}, \quad
        \hat{y} = \text{softmax}(W_s \cdot \text{Dropout}(\mathbf{h}_{\text{sent}}) + b_s)
    \end{equation}
    \item \textbf{ABSA}: Sử dụng biểu diễn token \texttt{[CLS]}:
    \begin{equation}
        \hat{Y}_{\text{absa}} = \text{reshape}\bigl(W_a \cdot \text{Dropout}(\mathbf{h}_{\text{CLS}}) + b_a\bigr) \in \mathbb{R}^{A \times P}
    \end{equation}
    với $A = 10$ aspects, $P = 3$ polarities. Loss là CrossEntropy trung bình qua tất cả aspects.
\end{itemize}

% -------------------------------------------------------------------
\subsection{Kiến trúc hệ thống}
% -------------------------------------------------------------------

\begin{figure}[H]
\centering
\begin{tikzpicture}[
    node distance=0.6cm and 1.2cm,
    block/.style={rectangle, draw=blue!70, fill=blue!8, rounded corners=4pt,
                  minimum width=2.8cm, minimum height=0.8cm, align=center,
                  font=\small},
    head/.style={rectangle, draw=orange!70, fill=orange!10, rounded corners=4pt,
                 minimum width=2.4cm, minimum height=0.7cm, align=center,
                 font=\small},
    data/.style={rectangle, draw=green!60!black, fill=green!8, rounded corners=4pt,
                 minimum width=2.2cm, minimum height=0.65cm, align=center,
                 font=\footnotesize},
    loss/.style={ellipse, draw=red!70, fill=red!8, minimum width=1.4cm,
                 minimum height=0.55cm, align=center, font=\footnotesize},
    arr/.style={-{Stealth[length=5pt]}, thick, draw=gray!70},
]

% Input
\node[block, minimum width=6cm] (input) {Input: \texttt{input\_ids}, \texttt{attention\_mask}};

% Backbone
\node[block, below=0.7cm of input, minimum width=6cm, fill=blue!15,
      draw=blue!80, font=\small\bfseries] (backbone)
      {PhoBERT Backbone (12 Transformer layers)\\$\mathbf{H} \in \mathbb{R}^{B \times L \times 768}$};

\draw[arr] (input) -- (backbone);

% Three heads
\node[head, below left=1cm and 1.5cm of backbone] (ner)
      {\textbf{NER Head}\\Linear + CRF};
\node[head, below=1cm of backbone] (sent)
      {\textbf{Sentiment Head}\\MeanPool + Linear};
\node[head, below right=1cm and 1.5cm of backbone] (absa)
      {\textbf{ABSA Head}\\CLS + Linear};

\draw[arr] (backbone.south) -- ++(0,-0.3) -| (ner.north);
\draw[arr] (backbone.south) -- (sent.north);
\draw[arr] (backbone.south) -- ++(0,-0.3) -| (absa.north);

% Outputs
\node[loss, below=0.6cm of ner]  (l1) {CRF NLL};
\node[loss, below=0.6cm of sent] (l2) {CrossEntropy};
\node[loss, below=0.6cm of absa] (l3) {CrossEntropy};

\draw[arr] (ner)  -- (l1);
\draw[arr] (sent) -- (l2);
\draw[arr] (absa) -- (l3);

% Output labels
\node[below=0.3cm of l1, font=\footnotesize\itshape, text=blue!70] {$[B,L,9]$ IOB2};
\node[below=0.3cm of l2, font=\footnotesize\itshape, text=blue!70] {$[B,3]$};
\node[below=0.3cm of l3, font=\footnotesize\itshape, text=blue!70] {$[B,10,3]$};

% Datasets
\node[data, above left=0.5cm and 0.6cm of input]  (d1) {VLSP NER};
\node[data, above=0.5cm of input]                  (d2) {UIT-VSFC};
\node[data, above right=0.5cm and 0.6cm of input]  (d3) {UIT-ViSFD};

\draw[arr, dashed, green!50!black] (d1.south) -- ++(0,-0.15) -| ([xshift=-1cm]input.north);
\draw[arr, dashed, green!50!black] (d2.south) -- (input.north);
\draw[arr, dashed, green!50!black] (d3.south) -- ++(0,-0.15) -| ([xshift=1cm]input.north);

\end{tikzpicture}
\caption{Kiến trúc MultiTaskPhoBERT: một backbone chia sẻ, ba head chuyên biệt.}
\label{fig:architecture}
\end{figure}

% -------------------------------------------------------------------
\subsection{Dữ liệu}
% -------------------------------------------------------------------

\begin{table}[H]
\centering
\caption{Thống kê ba bộ dữ liệu sử dụng.}
\label{tab:datasets}
\begin{tabular}{llccc}
\toprule
\textbf{Dataset} & \textbf{Task} & \textbf{Train} & \textbf{Val} & \textbf{Test} \\
\midrule
UIT-VSFC  & Sentiment (3 lớp)      & 11\,426 & 1\,538 & 3\,166 \\
UIT-ViSFD & ABSA (10 aspect $\times$ 3 pol.) & 8\,002 & 1\,001 & 1\,113 \\
VLSP 2016 & NER (IOB2: PER, ORG, LOC, MISC) & 14\,861 & 2\,000 & 2\,736 \\
\bottomrule
\end{tabular}
\end{table}

Tất cả dữ liệu được chuyển về một \texttt{DatasetDict} thống nhất
với schema chung \texttt{(text, task, sentiment\_label, aspect\_labels, ner\_tags)}.
Hàm \texttt{collate\_fn} tokenize bằng PhoBERT tokenizer, căn chỉnh nhãn sub-word
cho NER (chỉ token đầu tiên của mỗi từ nhận nhãn thật, các sub-word còn lại nhận \texttt{PAD\_LABEL\_ID = -100}).

% -------------------------------------------------------------------
\subsection{Chiến lược huấn luyện}
% -------------------------------------------------------------------

\begin{table}[H]
\centering
\caption{Các siêu tham số huấn luyện.}
\label{tab:hyperparams}
\begin{tabular}{ll}
\toprule
\textbf{Tham số} & \textbf{Giá trị} \\
\midrule
Backbone learning rate  & $2 \times 10^{-5}$ \\
Head learning rate      & $1 \times 10^{-3}$ \\
Optimizer               & AdamW (weight decay = 0.01) \\
Scheduler               & Linear warmup 10\% $\to$ Cosine decay \\
Gradient clipping       & max\_norm = 1.0 \\
Mixed precision         & AMP (\texttt{torch.amp.autocast}) \\
Batch size              & 16 \\
Max sequence length     & 256 \\
Epochs                  & 5 \\
Dropout (heads)         & 0.3 \\
Task sampling ratio     & NER 30\% / Sentiment 40\% / ABSA 30\% \\
\bottomrule
\end{tabular}
\end{table}

\paragraph{Alternating batch sampling.}
Thay vì huấn luyện tuần tự từng task,
mỗi bước lấy một batch từ task theo tỉ lệ cố định.
Cụ thể, trong mỗi epoch gồm $S$ steps tổng cộng
($S = \sum_t |\text{DataLoader}_t|$),
ta xây dựng lịch lặp (schedule) dạng round-robin
với tỉ lệ NER:Sentiment:ABSA = 3:4:3.

\paragraph{Class weights.}
Do phân phối nhãn mất cân bằng (đặc biệt lớp Neutral trong Sentiment chỉ chiếm $\sim$5\%),
ta tính inverse-frequency weights:
\begin{equation}
    w_c = \frac{N_{\text{total}}}{C \cdot N_c}
\end{equation}
và truyền vào CrossEntropy loss cho Sentiment và ABSA.

% -------------------------------------------------------------------
\subsection{Minh hoạ luồng xử lý}
% -------------------------------------------------------------------

\begin{figure}[H]
\centering
\begin{tikzpicture}[
    node distance=0.45cm,
    step/.style={rectangle, draw=blue!60, fill=blue!5, rounded corners=3pt,
                 minimum width=11cm, minimum height=0.65cm, align=left,
                 font=\footnotesize},
    arr/.style={-{Stealth[length=4pt]}, thick, draw=gray!60},
]

\node[step] (s1) {1. \textbf{Load data}: Đọc UIT-VSFC, UIT-ViSFD, VLSP-NER $\to$ \texttt{DatasetDict}};
\node[step, below=of s1] (s2) {2. \textbf{Tokenize}: PhoBERT BPE tokenizer, padding/truncation $\to$ \texttt{max\_length=256}};
\node[step, below=of s2] (s3) {3. \textbf{Build loaders}: Filter theo task, tạo DataLoader cho train/val mỗi task};
\node[step, below=of s3] (s4) {4. \textbf{Train}: Alternating batch $\to$ Forward backbone $\to$ Task head $\to$ Loss $\to$ Backward};
\node[step, below=of s4] (s5) {5. \textbf{Validate}: Mỗi cuối epoch, đánh giá F1 trên val set cho từng task};
\node[step, below=of s5] (s6) {6. \textbf{Checkpoint}: Lưu best model per-task khi val F1 cải thiện};
\node[step, below=of s6] (s7) {7. \textbf{Evaluate}: Load best checkpoint $\to$ Inference trên test set $\to$ Report};

\foreach \i/\j in {s1/s2, s2/s3, s3/s4, s4/s5, s5/s6, s6/s7} {
    \draw[arr] (\i) -- (\j);
}

\end{tikzpicture}
\caption{Pipeline xử lý end-to-end của Câu 1.}
\label{fig:pipeline}
\end{figure}

% -------------------------------------------------------------------
\subsection{Kết quả thực nghiệm}
% -------------------------------------------------------------------

Kết quả đánh giá trên tập test với checkpoint tốt nhất cho từng task:

\subsubsection{NER — VLSP 2016}

\begin{table}[H]
\centering
\caption{Kết quả NER entity-level trên test set (seqeval).}
\label{tab:ner_results}
\begin{tabular}{lcccc}
\toprule
\textbf{Entity} & \textbf{Precision} & \textbf{Recall} & \textbf{F1} & \textbf{Support} \\
\midrule
PER  & 0.9144 & 0.9380 & 0.9260 & 3\,884 \\
ORG  & 0.8887 & 0.8539 & 0.8710 & 3\,704 \\
LOC  & 0.8905 & 0.9193 & 0.9047 & 3\,717 \\
\midrule
\textbf{Macro avg} & \textbf{0.8979} & \textbf{0.9037} & \textbf{0.9006} & 11\,305 \\
\bottomrule
\end{tabular}
\end{table}

Mô hình đạt \textbf{Macro F1 = 0.9006} cho NER, vượt xa baseline all-O (F1 = 0.0).
Entity PER đạt kết quả tốt nhất (F1 = 0.926) do tên người thường có pattern rõ ràng trong tiếng Việt.
ORG là entity khó nhất (F1 = 0.871) do ranh giới tên tổ chức thường mơ hồ.

\subsubsection{Sentiment Analysis — UIT-VSFC}

\begin{table}[H]
\centering
\caption{Kết quả Sentiment Analysis trên test set.}
\label{tab:sent_results}
\begin{tabular}{lcccc}
\toprule
\textbf{Class} & \textbf{Precision} & \textbf{Recall} & \textbf{F1} & \textbf{Support} \\
\midrule
Negative & 0.9402 & 0.9595 & 0.9498 & 1\,409 \\
Neutral  & 0.6885 & 0.5030 & 0.5813 & 167 \\
Positive & 0.9477 & 0.9572 & 0.9524 & 1\,590 \\
\midrule
\multicolumn{2}{l}{\textbf{Macro F1}} & \multicolumn{2}{c}{\textbf{0.8278}} & \\
\multicolumn{2}{l}{\textbf{Weighted F1}} & \multicolumn{2}{c}{\textbf{0.9317}} & \\
\multicolumn{2}{l}{\textbf{Accuracy}} & \multicolumn{2}{c}{\textbf{0.9343}} & \\
\bottomrule
\end{tabular}
\end{table}

Accuracy đạt \textbf{93.4\%} (baseline majority-class: 50.2\%), cải thiện +43.2 điểm phần trăm.
Lớp Neutral có F1 thấp nhất (0.5813) do chỉ chiếm $\sim$5\% dữ liệu test,
gây khó khăn cho mô hình dù đã áp dụng class weights.

\subsubsection{ABSA — UIT-ViSFD}

\begin{table}[H]
\centering
\caption{Kết quả ABSA per-aspect trên test set.}
\label{tab:absa_results}
\begin{tabular}{lccc}
\toprule
\textbf{Aspect} & \textbf{F1 Macro} & \textbf{F1 Micro} & \textbf{Accuracy} \\
\midrule
SCREEN      & 0.8344 & 0.9438 & 94.4\% \\
CAMERA      & 0.9072 & 0.9487 & 94.9\% \\
FEATURES    & 0.8335 & 0.8840 & 88.4\% \\
BATTERY     & 0.9070 & 0.9209 & 92.1\% \\
PERFORMANCE & 0.8284 & 0.8323 & 83.2\% \\
STORAGE     & 0.7112 & 0.9915 & 99.2\% \\
DESIGN      & 0.8216 & 0.9353 & 93.5\% \\
PRICE       & 0.8188 & 0.9402 & 94.0\% \\
GENERAL     & 0.7994 & 0.8242 & 82.4\% \\
SER\&ACC    & 0.8265 & 0.9056 & 90.6\% \\
\midrule
\multicolumn{2}{l}{\textbf{Overall Micro F1}} & \multicolumn{2}{c}{\textbf{0.9126}} \\
\multicolumn{2}{l}{\textbf{Overall Macro F1}} & \multicolumn{2}{c}{\textbf{0.8661}} \\
\bottomrule
\end{tabular}
\end{table}

ABSA đạt \textbf{Overall Micro F1 = 0.9126} so với baseline majority-polarity (0.7345), cải thiện +17.8 điểm.
Aspect CAMERA đạt F1 macro cao nhất (0.907), trong khi STORAGE có F1 macro thấp nhất (0.711)
do phân phối cực kỳ lệch (hầu hết là neutral).

% -------------------------------------------------------------------
\subsection{So sánh với Baseline}
% -------------------------------------------------------------------

\begin{table}[H]
\centering
\caption{So sánh tổng hợp MultiTaskPhoBERT vs.\ các baseline đơn giản.}
\label{tab:comparison}
\begin{tabular}{lccc}
\toprule
\textbf{Task} & \textbf{Baseline} & \textbf{MultiTaskPhoBERT} & \textbf{$\Delta$} \\
\midrule
NER (Macro F1)         & 0.0000 (all-O) & \textbf{0.9006} & +0.9006 \\
Sentiment (Macro F1)   & 0.5022 (majority) & \textbf{0.8278} & +0.3256 \\
ABSA (Micro F1)        & 0.7345 (majority pol.) & \textbf{0.9126} & +0.1781 \\
\bottomrule
\end{tabular}
\end{table}

% -------------------------------------------------------------------
\subsection{Thông tin mô hình}
% -------------------------------------------------------------------

\begin{table}[H]
\centering
\caption{Thông tin kỹ thuật mô hình.}
\label{tab:model_info}
\begin{tabular}{ll}
\toprule
\textbf{Thuộc tính} & \textbf{Giá trị} \\
\midrule
Backbone            & \texttt{vinai/phobert-base} (RoBERTa, 12 layers) \\
Hidden size         & 768 \\
Model size          & 515.11\,MB \\
NER labels          & 9 (O, B/I-PER, B/I-ORG, B/I-LOC, B/I-MISC) \\
Sentiment labels    & 3 (Negative, Neutral, Positive) \\
ABSA output         & $10 \times 3 = 30$ logits \\
Latency (NER)       & $\sim$12.7\,ms/sample \\
Latency (Sent/ABSA) & $\sim$0.5\,ms/sample \\
\bottomrule
\end{tabular}
\end{table}

% -------------------------------------------------------------------
\subsection{Phân tích và thảo luận}
% -------------------------------------------------------------------

\paragraph{Ưu điểm của Multi-Task Learning.}
\begin{itemize}
    \item Chia sẻ backbone giúp mô hình học biểu diễn ngôn ngữ phong phú hơn
          từ nhiều tác vụ liên quan (NER cung cấp thông tin cú pháp, Sentiment cung cấp ngữ nghĩa).
    \item Tiết kiệm bộ nhớ: chỉ cần lưu 1 backbone ($\sim$500\,MB) thay vì 3 mô hình riêng.
    \item Regularization hiệu quả: mô hình phải đồng thời giải tốt 3 bài toán,
          giảm nguy cơ overfitting cho từng task.
\end{itemize}

\paragraph{Hạn chế.}
\begin{itemize}
    \item \textbf{Class imbalance}: Lớp Neutral trong Sentiment và một số aspect hiếm trong ABSA
          vẫn có F1 thấp dù đã dùng class weights.
    \item \textbf{MISC entity}: Không có support trong test set VLSP 2016, nên không đánh giá được.
    \item \textbf{Task interference}: Trong một số trường hợp, gradient từ task này
          có thể ``can thiệp'' tiêu cực đến task khác (negative transfer).
\end{itemize}

\paragraph{Hướng cải tiến.}
\begin{itemize}
    \item Thêm focal loss hoặc oversampling cho lớp thiểu số.
    \item Sử dụng task-specific adapters (LoRA/Adapter layers) thay vì fine-tune toàn bộ backbone.
    \item Thử nghiệm với PhoBERT-large ($H=1024$, 24 layers) cho kết quả cao hơn.
    \item Áp dụng gradient surgery hoặc PCGrad để giảm task interference.
\end{itemize}
